\section*{Chapter \ref{ch:g12:growth-and-production-function} --- Growth and the Production Function}
% figdata: figdata/grade-12/04-growth-and-production-function/out/ (cross, accounts, labsh, bren, bren_key); Brenland invented.

\begin{solution}{exo:g12:growth-and-production-function:1}
$64^{1/3} = 4$ and $8^{2/3} = 4$, so $Y = 2 \times 4 \times 4 = 32$.
\end{solution}

\begin{solution}{exo:g12:growth-and-production-function:2}
3 and 5.
\end{solution}

\begin{solution}{exo:g12:growth-and-production-function:3}
Capital: $461\,000/73\,000 \approx 6.3$. Output: $139\,000/19\,000 \approx 7.3$ (7.17 with the unrounded figures, as \cref{exo:g12:growth-and-production-function:14} uses).
\end{solution}

\begin{solution}{exo:g12:growth-and-production-function:4}
$1.5^{1/20} - 1 \approx 2.05\,\%$ a year.
\end{solution}

\begin{solution}{exo:g12:growth-and-production-function:5}
$2^{1/35} - 1 \approx 2.0\,\%$ a year.
\end{solution}

\begin{solution}{exo:g12:growth-and-production-function:6}
$F(\lambda K, \lambda L) = \lambda^{0.3 + 0.7}F = \lambda F$: constant returns. $G(\lambda K, \lambda L) = \lambda^{1.1}G$: increasing returns.
\end{solution}

\begin{solution}{exo:g12:growth-and-production-function:7}
$F(\lambda K, \lambda L) = \lambda F$: constant returns. $G(\lambda K, \lambda L) = \sqrt{\lambda}\,G$: decreasing returns.
\end{solution}

\begin{solution}{exo:g12:growth-and-production-function:8}
$3 - 0.3 \times 4 = 1.8\,\%$ a year.
\end{solution}

\begin{solution}{exo:g12:growth-and-production-function:9}
$4.55 - 2.50 = 2.05$ points (2.06 in the chart, from unrounded figures), about 45\,\% of growth.
\end{solution}

\begin{solution}{exo:g12:growth-and-production-function:10}
$2^{1/3} - 1 \approx 26\,\%$ each time: the second doubling adds 26\,\% of a larger $y$, but it needs twice as much extra capital as the first.
\end{solution}

\begin{solution}{exo:g12:growth-and-production-function:11}
United States: 0.628 in 1950, 0.637 in 2000, 0.568 in 2023. Korea: 0.422 in 1970, 0.546 in 2023.
\end{solution}

\begin{solution}{exo:g12:growth-and-production-function:12}
The American share fell from 0.637 in 2000 to 0.568 in 2023, and a study finds a decline since the early 1980s in most countries and industries. Under Cobb--Douglas the
share would stay at $1 - \alpha$ whatever happened to capital per worker. The evidence that cheaper investment goods explain about half the fall suggests that capital and
labour substitute more easily than the Cobb--Douglas form allows. (Korea's rising share shows that the movement is not the same everywhere.)
\end{solution}

\begin{solution}{exo:g12:growth-and-production-function:13}
$g_y = g_A + \alpha g_k + (1 - \alpha)g_h$. So $g_A = 2 - 2.5/3 - (2/3)(0.6) \approx 2 - 0.83 - 0.40 = 0.77$. Leaving out $h$, the residual is $2 - 0.83 = 1.17$: schooling
accounts for 0.40 points that would otherwise be counted as productivity.
\end{solution}

\begin{solution}{exo:g12:growth-and-production-function:14}
$6.32^{1/3} \approx 1.85$. To reach 7.17, $A$ must be $7.17/1.85 \approx 3.9$ times higher in the United States: capital explains a small part of the gap.
\end{solution}

\begin{solution}{exo:g12:growth-and-production-function:15}
$2^{1.1} \approx 2.14$. What the last unit of each input adds is $0.5Y/K$ and $0.6Y/L$, so paying them costs $0.5Y + 0.6Y = 1.1Y > Y$. Firms cannot pay every input its
competitive price; with increasing returns, large firms have lower costs per unit and markets tend away from \omterm{def:g11:perfect-competition:perfect-competition}{perfect competition}, unless the extra returns arise outside
the firm (for instance from ideas that spread).
\end{solution}

\begin{solution}{pb:g12:growth-and-production-function:1}
\textbf{1.} $34\,200/77 \times 100 \approx 44\,416$ million.

\textbf{2.} $K_{2000} = 3 \times 197\,403 = 592\,209$. $K_{2001} = 0.95 \times 592\,209 + 44\,416 \approx 607\,015$ million (607\,014 in the script's table, which uses
unrounded investment).

\textbf{3.} Machines and buildings wear out or become obsolete; without subtracting that, the stock would count capital that no longer produces.

\textbf{4.} $592\,209/3760 \approx 157.5$ and $953\,616/4366 \approx 218.4$ thousand dalens per worker.

\textbf{5.} $197\,403/3760 \approx 52.5$ and $292\,390/4366 \approx 66.97$ thousand dalens.

\textbf{6.} $\ln(66.97/52.5)/25 \approx 0.974\,\%$ a year.

\textbf{7.} $\ln(218.4/157.5)/25 \approx 1.308\,\%$ a year.

\textbf{8.} $0.35 \times 1.308 \approx 0.458$ points.

\textbf{9.} $0.974 - 0.458 \approx 0.516\,\%$ a year.

\textbf{10.} $0.516/0.974 \approx 53\,\%$.

\textbf{11.} $(1 - 0.35) \times 393\,264 \approx 255\,600$ million dalens.

\textbf{12.} $0.35 \times 292\,390/953\,616 \approx 0.107$: about 10.7\,\% a year, or 5.7\,\% after depreciation of 5\,\%.

\textbf{13.} $2^{0.35} - 1 \approx 27\,\%$.

\textbf{14.} $0.974/0.35 \approx 2.78\,\%$ a year, about twice as fast as it did.

\textbf{15.} Each addition to capital per worker raises output by less than the one before, so ever larger investment is needed for the same growth; without a rising $A$,
growth of output per worker fades.

\textbf{16.} A smaller starting stock makes each year's investment a larger addition to it, so capital per worker grows faster over the period; more of output growth is then
attributed to capital, and less is left for the residual.

\textbf{17.} $0.95^{25} \approx 0.28$: 28\,\% of the starting stock remains, about $0.28 \times 592\,209 \approx 164\,000$ of the 953\,616. Over a longer run the starting
stock decays further and the stock is made almost entirely of measured investment.

\textbf{18.} Errors in measuring output, capital or labour; the skills and schooling of workers; changes in how well resources are allocated among firms; the effects of
hours worked, of land and of natural resources left out of the function.

\textbf{19.} Over a few years, the residual mostly moves with the business cycle, as machines and workers are used more or less intensively; over sixty years those
swings average out.

\textbf{20.} \omterm{def:g12:growth-and-production-function:tfp}{Total factor productivity} accounts for about 53\,\% of Brenland's growth of output per worker from 2000 to 2025 (0.52 of 0.97 points a year).
\end{solution}
